\begin{conclusion}[有限域点计数的谱解释：仿射谱曲线点数与 Frobenius 迹的严格对应]\label{con:xi-terminal-zm-elliptic-finite-field-point-count}
沿用结论 \ref{con:xi-terminal-zm-elliptic-normalization} 的谱四次 \(\Pi(\lambda,y)\) 与极小椭圆模型
\[
E_0:\ v^{2}+v=x^{3}-x.
\]
取素数 \(p\neq 37\)，并记在 \(\FF_p\) 上的仿射解集
\[
C_p(\FF_p):=\{(\lambda,y)\in\FF_p^{2}:\ \Pi(\lambda,y)=0\},
\qquad
E_{0,p}(\FF_p):=\{(x,v)\in\FF_p^{2}:\ v^{2}+v=x^{3}-x\}\cup\{O\}.
\]
则由整系数多项式互逆代换
\[
(x,v)\longmapsto(\lambda,y)=(x,\ v+x^{2}),
\qquad
(\lambda,y)\longmapsto(x,v)=(\lambda,\ y-\lambda^{2})
\]
在 \(\FF_p\) 上诱导双射
\[
E_{0,p}(\FF_p)\setminus\{O\}\ \longrightarrow\ C_p(\FF_p),
\]
从而
\[
\#C_p(\FF_p)=\#E_{0,p}(\FF_p)-1.
\]
令
\(
a_p:=p+1-\#E_{0,p}(\FF_p)
\)
为 \(E_0\) 在良素 \(p\neq 37\) 处的 Frobenius 迹，则等价地有
\[
\#C_p(\FF_p)=p-a_p.
\]
并且对奇素数 \(p\neq 37\) 成立二次特征和公式
\[
a_p=-\sum_{x\in\FF_p}\left(\frac{4x^{3}-4x+1}{p}\right),
\]
其中 \(\bigl(\tfrac{\cdot}{p}\bigr)\) 为 Legendre 符号（约定 \(\bigl(\tfrac{0}{p}\bigr)=0\)）。
\end{conclusion}
\begin{proof}[证明要点]
变换 \((\lambda,y)\leftrightarrow(x,v)\) 由结论 \ref{con:xi-terminal-zm-elliptic-normalization} 给出，且在 \(\ZZ\) 上为多项式互逆；故对任意素数 \(p\) 均在 \(\FF_p\) 上诱导仿射点集的双射。
由于 Weierstrass 模型 \(E_{0,p}\) 的射影完备化仅多出唯一无穷远点 \(O\)，得 \(\#C_p(\FF_p)=\#E_{0,p}(\FF_p)-1\)，从而 \(\#C_p(\FF_p)=p-a_p\)。
\par
设 \(p\neq 2,37\)。
对每个 \(x\in\FF_p\)，方程 \(v^{2}+v=x^{3}-x\) 的解数等于
\[
1+\left(\frac{1+4(x^{3}-x)}{p}\right)=1+\left(\frac{4x^{3}-4x+1}{p}\right),
\]
因此
\[
\#E_{0,p}(\FF_p)
=1+\sum_{x\in\FF_p}\left(1+\left(\frac{4x^{3}-4x+1}{p}\right)\right)
=p+1+\sum_{x\in\FF_p}\left(\frac{4x^{3}-4x+1}{p}\right).
\]
代入 \(a_p=p+1-\#E_{0,p}(\FF_p)\) 得所示公式。证毕。
\end{proof}

\begin{proposition}[有限域纤维计数与 Frobenius 迹的均值恒等式]\label{prop:xi-terminal-zm-elliptic-fiber-count-mean-trace}
沿用结论 \ref{con:xi-terminal-zm-elliptic-finite-field-point-count} 的记号，并取良素 \(p\neq 37\)。
对每个 \(y\in\FF_p\) 定义纤维计数
\[
N_p(y):=\card{\{\lambda\in\FF_p:\ \Pi(\lambda,y)=0\}}.
\]
则成立：
\[
\#E_{0,p}(\FF_p)=1+\sum_{y\in\FF_p}N_p(y),
\qquad
\sum_{y\in\FF_p}\bigl(N_p(y)-1\bigr)=-a_p.
\]
\leanverified{paper\_xi\_terminal\_zm\_elliptic\_fiber\_count\_mean\_trace}
\end{proposition}

\begin{proof}
按定义，
\(
C_p(\FF_p)=\{(\lambda,y)\in\FF_p^2:\ \Pi(\lambda,y)=0\}
\)
的点数可按 \(y\)-纤维分解为
\[
\#C_p(\FF_p)=\sum_{y\in\FF_p}N_p(y).
\]
而结论 \ref{con:xi-terminal-zm-elliptic-finite-field-point-count} 已给出 \(\#C_p(\FF_p)=\#E_{0,p}(\FF_p)-1\)，故得第一式。
再由同一结论的 \(\#C_p(\FF_p)=p-a_p\) 得
\[
\sum_{y\in\FF_p}\bigl(N_p(y)-1\bigr)
=\#C_p(\FF_p)-p
=(p-a_p)-p
=-a_p.
\]
证毕。
\end{proof}

\begin{corollary}[Hecke 迹的随机纤维估计：无偏性与方差公式]\label{cor:xi-terminal-zm-elliptic-mc-trace-estimator}
\leanverified{paper\_xi\_terminal\_zm\_elliptic\_mc\_trace\_estimator}
在命题 \ref{prop:xi-terminal-zm-elliptic-fiber-count-mean-trace} 的设定下，
令 \(Y_1,\dots,Y_n\) 为 \(\FF_p\) 上独立均匀样本，并令
\[
Z_i:=N_p(Y_i)-1,
\qquad
\widehat a_p:=-p\cdot \frac{1}{n}\sum_{i=1}^{n} Z_i.
\]
则 \(\widehat a_p\) 为 \(a_p\) 的无偏估计，且
\[
\EE[\widehat a_p]=a_p,
\qquad
\Var(\widehat a_p)=\frac{p^{2}}{n}\Var(Z_1).
\]
\end{corollary}

\begin{proof}
由命题 \ref{prop:xi-terminal-zm-elliptic-fiber-count-mean-trace}，
\[
\EE[Z_1]
=\frac{1}{p}\sum_{y\in\FF_p}\bigl(N_p(y)-1\bigr)
=-\frac{a_p}{p},
\]
故 \(\EE[\widehat a_p]=a_p\)。
方差式由独立同分布与缩放恒等式直接推出：
\(
\Var(\widehat a_p)
=p^2\Var(\frac1n\sum_i Z_i)
=p^2\cdot \frac{1}{n}\Var(Z_1)
\)。
证毕。
\end{proof}
